% !TEX root = ../main.tex
\chapter{Objetivos}
\section{Objetivo general}

Desarrollar un modelo semi-supervisado de redes neuronales convolucionales equivariantes para la segmentación y clasificación de imágenes histopatológicas.


\section{Objetivos específicos}

\begin{itemize}
    \item Seleccionar un conjunto de datos de imágenes histopatológicas anotadas que hayan sido recopiladas para la clasificación y segmentación de estructuras microscópicas.

    \item Formular e implementar un modelo semi-supervisado basado en una arquitectura de red neuronal convolucional equivariante para la segmentación y clasificación de imágenes histopatológicas.

    \item Evaluar el desempeño del modelo de redes neuronales convolucionales equivariantes sobre un conjunto de validación de imágenes histopatológicas a través de métricas para la clasificación y segmentación de imágenes.

\end{itemize}

Para conservar la formulación aprobada en Tesis I y, al mismo tiempo, precisar
su correspondencia con el trabajo ejecutado, el término \emph{modelo
semisupervisado} designa aquí un aprendizaje no supervisado de representaciones
seguido por evaluaciones supervisadas sobre codificadores congelados. La
clasificación se realizó a nivel de WSI y la segmentación se operacionalizó como
clasificación espacial de celdas sobre una cuadrícula de parches. Esta última no
es exhaustiva porque las máscaras disponibles no anotan toda la extensión de las
láminas.
